\chapter{Appendinx: Draft}

\subsection{what and why}
what and why CDN

\subsection{siby2023quic}
paper~\cite{siby2023quic}

\textbf{\textit{QUIC.}} QUIC is a connection-oriented protocol built on top of UDP
that aims to provide low-latency, multiplexed, secure communication
with less head-of-line blocking and faster connection migration [7].
QUIC was standardized in May 2021 by the IETF. QUIC is the trans-
port protocol for HTTP/3. Adoption of QUIC and HTTP/3 has been
rising (as of June 2023, they are used by 26\% of the top 10 million
websites [17]). Of particular relevance for our work is the QUIC
PADDING Frame. The IETF QUIC standard describes it as a frame
with no semantic value, that can be used to increase packets size and
to provide protection against traffic analysis [8].

Paper did simplified test, for general public, which included testing different attacks and defenses without using VPN, which negatively impacted defence, and proved that Padding alone is not sufficient.

They tried to identify how QUIC works and how padding can help defend against traffic analysis. First, they measured the baseline without any protection and got over 90\% accuracy using a Random Forest classifier. Then they added padding, but it barely helped, as they still got around 80\%. After that, they tried adding fake packets, which brought the accuracy down to around 30\%. I guess that's an improvement, but it's still far from ideal, considering random guessing would be below 1\%.

The biggest problem, though, is that they never measured the latency overhead, and the additional padding almost doubled the bandwidth usage.



\subsection{https://pulls.name/blog/2024-06-05-eval-first-daita-servers/}
paper~\cite{pulls2024daitaeval}

How we combat traffic analysis: this is how DAITA works.

\textbf{\textit{DAITA}} has been developed together with Computer Science at Karlstad University and uses three types of cover traffic to resist traffic analysis.
\begin{enumerate}
    \item  Constant Packet Sizes
The size of network packets can be particularly revealing, especially small packets, so DAITA makes all packets sent over the VPN the same constant size.
    \item  Random Background Traffic
 By unpredictably interspersing dummy packets into the traffic, DAITA masks the routine signals to and from your device. This makes it harder for observers to distinguish between meaningful activity and background noise.
    \item  Data Pattern Distortion
When visiting websites (or doing any other activity that causes significant traffic), DAITA modifies the traffic pattern by unpredictably sending cover traffic in both directions between client and VPN server. This distorts the recognizable pattern of a website visit, resisting accurate identification of the site. [from mullvad paper itself]
\end{enumerate}

\textbf{DAITA} reduces the success rate of both DF and RF attacks, with more prolonged traces (US servers) offering slightly less protection. The protection costs around 2-4 MiB of bandwidth overhead and 1-3 seconds of delay per website visit, depending on server (and website) location.

The defenses are padding-only defenses that create so much padding traffic that congestion occurs, implicitly introducing notable delays. In the future, defenses should explicitly create short, useful (for protection) delays and aim to account for congestion to minimize delays.


\textbf{negative result of blogs experiment}
The user experience of enabling DAITA when browsing with Mullvad Browser is less than changing the server location from Europe to the US. Recall that the Karlstad University website is hosted in Sweden, so the traffic crossed the Atlantic twice in our experiment (unlike most regular website visits most users do). Not great, but not terrible.

\subsection{paper2}
paper~\cite{dyer2012peekaboo}

\subsection{paper3}
paper~\cite{cai2014tamaraw}

\subsection{paper4}
paper~\cite{gong2020front}
